%------------------------------------------------------------------------------%
\documentclass[fleqn,utf8,aspectratio=169,12pt]{beamer}

\usepackage{gaal}

\title{Equações Lineares}

%------------------------------------------------------------------------------%
\begin{document}

\addtitle

%------------------------------------------------------------------------------%
\begin{frame}{Equações Lineares}
  Equação linear com duas incógnitas
  \[
    ax + by = 0
  \]

  \pause
  Equação linear com $n$ incógnitas
  \[
    a_1x_1 + a_2x_2 + \cdots + a_nx_n = 0
  \]

  \pause
  Possuem infinitas soluções
\end{frame}

%------------------------------------------------------------------------------%
\begin{frame}{Equação Linear Não Homogênea com Duas Incógnitas}
  \onslide<+->{Equação linear não homogênea com duas incógnitas
  \[
    ax + by = c
  \]}
  \pause
  \onslide<+->{Possui solução única se $a$ e $b$ não forem zero}
  \begin{align*}
    \onslide<+->{ax + by             & = c                \\}
    \onslide<+->{ax + by - \emph{ax} & = c -\emph{ax}     \\}
    \onslide<+->{by                  & = c - ax           \\}
    \onslide<+->{\frac{by}{\emph{b}} & = \frac{c - ax}{\emph{b}} & \text{desde que, $b$ não seja zero} \\}
    \onslide<+->{y                   & = \frac{c - ax}{b}}
  \end{align*}
\end{frame}

%------------------------------------------------------------------------------%
\begin{frame}{Sistema de 3 Equações com 3 Incógnitas}
  \systeme{
    110 x + 120 y + 130 z = 100,
    210 x + 220 y + 230 z = 200,
    310 x + 320 y + 330 z = 300
  }
\end{frame}

%------------------------------------------------------------------------------%
\begin{frame}{Sistema de $m$ Equações com $n$ Incógnitas}
  \begin{align*}
    a_{11} x_1 + a_{21} x_2 + \cdots + a_{n1}x_n & = b_1 \\
    a_{12} x_1 + a_{22} x_2 + \cdots + a_{n2}x_n & = b_2 \\
    a_{13} x_1 + a_{23} x_2 + \cdots + a_{n3}x_n & = b_3 \\
    a_{14} x_1 + a_{24} x_2 + \cdots + a_{n4}x_n & = b_4 \\
    & \vdots  \\
    a_{1m} x_1 + a_{2m} x_2 + \cdots + a_{nm}x_n & = b_m
  \end{align*}
\end{frame}

%------------------------------------------------------------------------------%
\begin{frame}{Representação Matricial}
  \[
    A = \begin{pmatrix}
      a_{11} & a_{21} & \cdots & a_{n1} \\
      a_{12} & a_{22} & \cdots & a_{n2} \\
      a_{13} & a_{23} & \cdots & a_{n3} \\
      a_{14} & a_{24} & \cdots & a_{n4} \\
      \vdots & \vdots & \vdots & \vdots \\
      a_{1m} & a_{2m} & \cdots & a_{nm}
    \end{pmatrix}
    \pause
    \qquad\qquad
    x = \begin{pmatrix}
      x_1    \\
      x_2    \\
      x_3    \\
      x_4    \\
      \vdots \\
      x_m
    \end{pmatrix}
    \pause
    \qquad\qquad
    b = \begin{pmatrix}
      b_1    \\
      b_2    \\
      b_3    \\
      b_4    \\
      \vdots \\
      b_m
    \end{pmatrix}
  \]

  \pause
  \bigskip
  \[
    Ax = b
  \]
\end{frame}

%------------------------------------------------------------------------------%
%\section{Lista Mínima}
%
%%------------------------------------------------------------------------------%
%\begin{frame}{Lista Mínima}
%  \begin{enumerate}
%    \item
%  \end{enumerate}
%
%  \vfill
%  Atenção: A prova é baseada no livro, não nas apresentações
%\end{frame}

%------------------------------------------------------------------------------%
\end{document}
%------------------------------------------------------------------------------%
